\section{Постановка технического задания для системы МОТ}\label{sec:task}

Таким образом, проведя анализ архитектуры системы управления автономным ТС, сенсоров, систем детекции, а также систем трекинга можно сформировать техническое задание на ВКРМ. Учитывая аппаратную платформу, на которой будет производится запуск системы детекции и отслеживания, разрабатываемая система должна удовлетворять следующим требованиям:

\begin{enumerate}
    \item Скорость работы на двух чипах Intel Myriad должна быть не менее 25 FPS, при разрешении входного видеопотока 640 на 640 пикселей.
    \item Показатели результатов обучения детектора:
     \begin{itemize}
        \item mAP\@0.5 > 0.8
        \item IDF1 > 0.78
        \item HOTA > 0.64
        \item FAF < 0.16
    \end{itemize}
    \item Средние показатели по всем тестировочным сценариям:
    \begin{itemize}
        \item MOTA > 0.81
        \item IDF1 > 0.78
        \item HOTA > 0.64
        \item FAF < 0.16
        \item IDS < 30/1000 кадров
    \end{itemize}
    \item Средние показатели по нагруженным сценариям (плохое освещение, плотный трафик, плохие погодные условия):
    \begin{itemize}
        \item MOTA > 0.76
        \item IDF1 > 0.75
        \item HOTA > 0.6
        \item FAF < 0.16
        \item IDS < 50/1000 кадров
    \end{itemize}
\end{enumerate}

Технические характеристики IntelMyriad представлены в приложении~\ref{tab:myriadx_specs}~\cite{Intel_MyriadX_ProductBrief_2018}~\cite{Oh_2017_MyriadX_AnandTech}.

Также, в рамках проекта было предложено разработать функцию для поиска наиболее важного объекта (НВО). Данная функция должна:

\begin{itemize}
    \item Анализировать скорость и положение окружающих ТС;
    \item Исключать статические объекты;
    \item Иметь возможность принимать параметры ТС для возможности интеграции на различные ТС;
    \item Принимать данные с CAN-шины ТС;
    \item В качестве выходных данных возвращать параметры НВО (ID, положение, скорость).
\end{itemize}